%
% Documento: Introdução
%
\chapter{Introdução}\label{chap:introducao}

A energia elétrica constitui um insumo essencial para o funcionamento da sociedade contemporânea, sustentando atividades produtivas, serviços públicos e aplicações críticas em áreas como saúde, comunicação e infraestrutura. Nos sistemas de distribuição, a continuidade do fornecimento e a confiabilidade da operação dependem da capacidade de identificar perturbações que, embora nem sempre apresentem correntes elevadas, podem comprometer a segurança e a integridade da rede.

Nesse contexto, as falhas de alta impedância (FAIs) são reconhecidas, desde os trabalhos clássicos da década de 1980, como um dos problemas mais persistentes e difíceis enfrentados pela proteção de sistemas elétricos. \cite{1,2} Essas ocorrências foram descritas como situações em que um condutor energizado entra em contato com meios de alta resistência, como solo seco, asfalto ou outros materiais semelhantes, produzindo correntes insuficientes para a atuação confiável das proteções convencionais por sobrecorrente. \cite{3} Além disso, a literatura demonstra que as FAIs apresentam características complexas, com comportamento não linear, assimetria e assinaturas dependentes das condições de contato e do ambiente, o que dificulta sua detecção e localização. \cite{4,5}

Com o avanço das pesquisas a partir dos anos 2000, diferentes métodos passaram a ser propostos para melhorar a identificação dessas ocorrências, incluindo abordagens baseadas em componentes de alta frequência, harmônicos, inter-harmônicos e análise no domínio do tempo e da frequência. \cite{5,6,7,8} Esses estudos indicam que, embora haja progresso na modelagem e na detecção, o problema permanece relevante para a engenharia de proteção, sobretudo quando se consideram as variações reais do sistema e as dificuldades de implementação em campo. \cite{6,7}

Quando associadas à vegetação, as FAIs assumem importância ainda maior, pois podem prolongar a ocorrência de arco elétrico e aumentar o risco de ignição de materiais combustíveis. A literatura recente destaca que as falhas relacionadas à vegetação podem causar danos à rede, dificuldades operacionais e até incêndios, o que as torna particularmente críticas em áreas com arborização próxima às linhas aéreas. \cite{8,9}

Os dados australianos utilizados neste trabalho \cite{vic_bushfire_report} foram obtidos a partir do programa experimental Vegetation Conduction Ignition Test Program, desenvolvido no âmbito do Powerline Bushfire Safety Program do governo de Victoria, após os incêndios florestais de 2009. Esse programa foi criado para aprofundar o entendimento dos processos de ignição em falhas envolvendo vegetação e para estabelecer uma base de referência de assinaturas elétricas de falha, com o objetivo de apoiar o desenvolvimento de métodos mais eficientes de detecção. O conjunto de dados é especialmente relevante porque contém registros reais de testes controlados, com múltiplas espécies vegetais, diferentes condições experimentais e sinais elétricos capturados em alta resolução, o que o torna adequado para a formulação de um problema de classificação.

Diante do exposto, o presente estudo propõe utilizar esses registros para construir um problema de classificação supervisionada, no qual as medições elétricas associadas às ocorrências de falha serão analisadas a fim de distinguir classes relacionadas à presença ou ausência de contato vegetação–linha, ou ainda categorias derivadas do tipo de vegetação, severidade da ignição ou estágio do evento.